\begin{abstract}
The value of Planck's constant was obtained using the wavelengths of the first three emission lines of the hydrogen atom, 
determined from their diffraction angles. Similarly, the wavelength of an unknown element was measured in order to 
identify its spectral series and consequently determine which element it corresponds to.
\end{abstract}
\begin{IEEEkeywords}
Planck's constant, spectral series, diffraction angle, and wavelength.
\end{IEEEkeywords}

\selectlanguage{spanish}
\begin{abstract}
Se obtuvo el valor de la constante de planck por medio de las longitudes de onda de las primeras tres líneas de emisión del
átomo de hidrógeno conseguidas a partir de su ángulo de difracción. De manera análoga, se determinó la longitud de onda de un 
elemento desconocido con el fin de identificar su serie espectral y, en consecuencia, determinar a qué elemento corresponde.
\end{abstract}
\begin{IEEEkeywords}
Constante de Planck, serie espectral, ángulo de difracción y longitud de onda.
\end{IEEEkeywords}
